\chapter{多模态数据}
\label{lab:E20_multimodal_llm}

\noindent\textbf{配套文件：}\path{notebook/E20_multimodal_llm.ipynb}\par
\noindent\textbf{教材关联：}\bookxrefshort{ch:multimodal-alignment}。

\section*{实验目标}
验证视觉表示如何进入语言模型的序列与监督区域。

\section*{准备及定位}
先阅读本册实验总则，确认Notebook中的依赖与资产单元。原理计算、库接口迁移和真实模型训练可能具有不同资源要求，应分别记录设备、版本及运行范围。

源码定位可从下列函数开始；应结合前置数据与配置单元阅读。
\begin{itemize}
\item \texttt{my\_draw\_connector\_block}
\item \texttt{my\_draw\_connector\_arrow}
\end{itemize}

\section*{操作及观察}
\begin{enumerate}
\item 记录图像输入与patch token的形状，区分视觉分辨率与词元长度。
\item 对照Projector和Resampler的维度与长度变换。
\item 组装视觉和文本序列，逐项检查占位符、注意力可见性及标签掩码。
\item 比较冻结与解冻阶段的可训练模块，审查迁移到标准处理器时的协议变化。
\end{enumerate}

\section*{结果判读及提交}
提交形状流、掩码与参数范围。张量连接正确不证明视觉问答能力；模型图谱快照不能当作持续更新的排名。

提交时附上所执行单元范围、环境与制品身份、原始结果及失败说明。将未运行项目明确列出，不能用Notebook中已有输出替代本次执行证据。
